% Pursuing Stacks (cotes 134-2 to 134-8): archive cote 134-8, site page 17 (PDF page 18), typescript page 571.
% typescript. Draft compared: ps7.tex:785-840.
% AI-transcribed by Claude Opus 5.5 (claude-opus-5-5, adaptive thinking, effort high), Message Batches API, batch ps4 (msgbatch_012Aoc1NuskATPsyBeKSYK7J), 2026-11-12; not yet checked by eye.
\scanpage{134-8/17}
\mspage{571}

cohomological counterparts of smoothness and properness fit very neatly
into the homology and cohomology formalism,  and I played around enough
with them, \xout{xxxx}\inserted[typed]{last} spring as well as more than twenty years ago when
developping étale cohomology of schemes, \over{\unsure{s}}{f}o\hand{r} there \deleted{is indeed}\inserted{being} no doubt
left in my mind that these notions are relevant indeed.  However I
\xout{forgot} was rather rash yesterday,  while forgetting that the\hand{\unsure{se}}cohomologi
-cal versions of smoothness and properness are considerably weaker than
the usual notions.  Thus in the context of schemes over a ground field,
the product of any two schemes is cohomologically smooth over it's
factors  -  or equivalently,  any scheme over a field\inserted{k} is cohomological-%
ly smooth over k !  Similarly,  in the context of (Cat),  any \xout{category}
object in (Cat)\hand{,} namely any small category\hand{,} is both cohomologically
smooth and proper over the final object $e$  (as it is trivially \hand{``}bifi-%
bered\hand{''}over $e$ ).  On the other hand,  it is'nt reasonable\hand{,} of course to
expect any kind of Poincaré-like duality to hold for the cohomology
(with twisted coefficients say) of an arbitrary object A in \hand{(Cat).} To be
more specific,  it is easy to see that in many cases,  the functor
\[ Rf_{*}\ \text{\xout{(xxxxRf)}} : D^{+}(A^{\wedge}_{ab}) \longrightarrow D^{+}(Ab) \]
does not admit a right adjoint (which we would call $f^{!}$ ).  For instance
when A is discrete,  then $f^{!}$  exists (and may then be identified with
$f^{*}$ ) iff A is moreover \uline{finite}  -  a rather natural condition indeed,
when we keep in mind the topological significance of the usual notion
of properness !  This immediately brings to mind some further proper-%
ties besides base change properties,  which go with the intuitions
around properness  -  for instance we would expect for proper f the
functors $f_{*}$ and $Rf_{*}$ to commute to filtering direct limits, and the
same expectation goes with the assumption that $Rf_{*}$ should admit a
right adjoint.  This exactness property is not satisfied of course
when A is di\xout{\unsure{x}}screte infinite. \xout{Another}  We\inserted{now} may think (still in case of
a target category equal to $e$) to impose the drastic condition that the
category A is finite.  Such restriction however looks in some respects
too weak,  in others too strong.  Thus it will include categories de-%
fined by finite groups,\deleted{\unsure{f}} which goes against the rather natural expec-%
tation that properness + smoothness,  or any kind of Poincaré duality,
should go with \uline{finite cohomological dimension}.  On the other hand,
there are beautiful infinite groups (such as the fundamental group\xout{s}
of a  compact surface,  or of any other compact variety which is a
$K(\hand{$\pi$},1)$ space ...) which satisfy Poincaré duality.

\margin{17}
